\begin{mistake}{2026-07-17}{02}

\begin{problemcontent}
设函数 \(f(x)\) 在区间 \((a,b)\) 内可导，行列式 \(|A| = \begin{vmatrix} 1 & x_1 & f(x_1) \\ 1 & x_2 & f(x_2) \\ 1 & x_3 & f(x_3) \end{vmatrix}\)。证明：导函数 \(f'(x)\) 在 \((a,b)\) 内严格单调递增的充分必要条件是对 \((a,b)\) 内任意的 \(x_1 < x_2 < x_3\)，有 \(|A| > 0\)。
\end{problemcontent}

\begin{solutioncontent}
\textbf{第一步：行列式化简}
对 \(|A|\) 进行初等行变换（\(r_2 - r_1, r_3 - r_2\)），按第一列展开可得：
\[
\begin{aligned}
|A| &= \begin{vmatrix} 1 & x_1 & f(x_1) \\ 0 & x_2-x_1 & f(x_2)-f(x_1) \\ 0 & x_3-x_2 & f(x_3)-f(x_2) \end{vmatrix} \\
    &= (x_2-x_1)(x_3-x_2) \left[ \frac{f(x_3)-f(x_2)}{x_3-x_2} - \frac{f(x_2)-f(x_1)}{x_2-x_1} \right]
\end{aligned}
\]
因 \(x_1 < x_2 < x_3\)，故 \((x_2-x_1)(x_3-x_2) > 0\)。因此 \(|A| > 0\) 等价于：
\[ \frac{f(x_3)-f(x_2)}{x_3-x_2} > \frac{f(x_2)-f(x_1)}{x_2-x_1} \quad \text{(式1)} \]

\textbf{第二步：必要性 (\(\Rightarrow\))}
已知 \(f'(x)\) 严格单调递增。对 \([x_1, x_2]\) 和 \([x_2, x_3]\) 使用拉格朗日中值定理，存在 \(\xi_1 \in (x_1, x_2)\), \(\xi_2 \in (x_2, x_3)\)：
\[ \frac{f(x_2)-f(x_1)}{x_2-x_1} = f'(\xi_1), \quad \frac{f(x_3)-f(x_2)}{x_3-x_2} = f'(\xi_2) \]
因 \(\xi_1 < x_2 < \xi_2\) 且 \(f'(x)\) 严格递增，故 \(f'(\xi_1) < f'(\xi_2)\)，即（式1）成立，从而 \(|A| > 0\)。

\textbf{第三步：充分性 (\(\Leftarrow\))（利用凹函数几何性质法）}
已知对任意 \(x_1 < x_2 < x_3\)，（式1）恒成立。这表明 \(f(x)\) 的割线斜率严格递增，即 \(f(x)\) 为严格凹函数（开口向上）。
对于可导的严格凹函数，其图像始终严格位于切线上方。任取 \(y_1 < y_2\)，分别在 \(y_1, y_2\) 处作切线：
\[ f(y_2) > f(y_1) + f'(y_1)(y_2 - y_1) \]
\[ f(y_1) > f(y_2) + f'(y_2)(y_1 - y_2) \]
两式相加并化简得：
\[ 0 > [f'(y_1) - f'(y_2)](y_2 - y_1) \]
因 \(y_2 - y_1 > 0\)，故必有 \(f'(y_1) < f'(y_2)\)。充分性得证。
\end{solutioncontent}

\mistakeknowledge{
\begin{itemize}
  \item \textbf{核心方法}：利用初等行变换将行列式转化为差商之差；利用拉格朗日中值定理建立差商与导数的关系。
  \item \textbf{易错点}：充分性证明中，若直接用导数定义取极限，严格不等号会变成“\(\le\)”。必须引入中间点，或利用严格凹函数的切线性质（图像在切线严格上方）来避免取极限带来的等号陷阱。
  \item \textbf{关联知识}：割线斜率严格递增是严格凹函数（即下凸函数）的等价定义；严格凹函数满足 \(f(x) > f(x_0) + f'(x_0)(x-x_0)\) (\(x \neq x_0\))。
\end{itemize}}

\end{mistake}
